% \usepackage[notextcomp,nomath]{stix2}
\usepackage{tcolorbox}
\usepackage{amsfonts,amscd,amssymb,amsmath,amsthm,enumerate}
\numberwithin{equation}{section}
\usepackage{array,mathtools,mathrsfs,bm}
\usepackage{threeparttable}
\usepackage{booktabs,caption,longtable,multicol,multirow}
\usepackage{subcaption,lscape}
\usepackage{makecell}
\usepackage{graphicx}
\usepackage{enumerate}
\usepackage{enumitem}
\setlist[enumerate]{noitemsep, topsep=2pt}
\setlist[itemize]{noitemsep, topsep=2pt}
% \usepackage[linesnumbered,boxruled,lined,commentsnumbered]{algorithm2e}
\usepackage[colorlinks,citecolor=purple,linkcolor=blue]{hyperref}
% \usepackage[nameinlink]{cleveref}
\usepackage{cleveref}
% \usepackage[parfill]{parskip}
\usepackage[margin=1.5in]{geometry}
\usepackage[sort&compress,numbers]{natbib}
% \usepackage{authblk}
\usepackage{algorithm,algorithmic}
\usepackage{pifont}
\usepackage{appendix}
\usepackage{multirow}
\usepackage{tabularx}
\usepackage{makecell}
\usepackage{tikz}
\usepackage{pgfplots}
\usetikzlibrary{arrows}
\usetikzlibrary{calc}
\usetikzlibrary{arrows.meta}
\usetikzlibrary{backgrounds}
\usepgfplotslibrary{patchplots}
\usepgfplotslibrary{fillbetween}
% \usepackage[inline]{showlabels}
% \renewcommand{\showlabelfont}{\tiny\color{blue}}
%%%%%%%%%%%%%%%%%%%%%%%%%%%%%%%%%%%%%%%%%%%%%%%%
\usepackage{eqparbox}
\renewcommand\algorithmiccomment[1]{%
  \hfill\texttt{\#}\ \eqparbox{COMMENT}{\texttt{#1}}%
}
%%%%%%%%%%%%%%%%%%%%%%%%%%%%%%%%%%%%%%%%%%%%%%%%

%% tikz
\definecolor{ao(english)}{rgb}{0.0, 0.5, 0.0}
\definecolor{cadmiumgreen}{rgb}{0.0, 0.42, 0.24}
\definecolor{darkpastelgreen}{rgb}{0.01, 0.75, 0.24}
\let\subfigureautorefname=\figurename
\usepackage{adjustbox}
% Define block styles
\tikzset{%
  materia/.style={draw, fill=blue!20, text width=6.0em, text centered, minimum height=1.5em,drop shadow},
  etape/.style={materia, text width=8em, minimum width=10em, minimum height=3em, rounded corners, drop shadow},
  texto/.style={above, text width=6em, text centered},
  linepart/.style={draw, thick, color=black!50, -LaTeX, dashed},
  line/.style={draw, thick, color=black!50, -LaTeX},
  ur/.style={draw, text centered, minimum height=0.01em},
  back group/.style={fill=white!20,rounded corners, draw=black!50, dashed, inner xsep=15pt, inner ysep=10pt},
}
\pgfplotsset{%
  layers/standard/.define layer set={%
      background,axis background,axis grid,axis ticks,axis lines,axis tick labels,pre main,main,axis descriptions,axis foreground%
    }{
      grid style={/pgfplots/on layer=axis grid},%
      tick style={/pgfplots/on layer=axis ticks},%
      axis line style={/pgfplots/on layer=axis lines},%
      label style={/pgfplots/on layer=axis descriptions},%
      legend style={/pgfplots/on layer=axis descriptions},%
      title style={/pgfplots/on layer=axis descriptions},%
      colorbar style={/pgfplots/on layer=axis descriptions},%
      ticklabel style={/pgfplots/on layer=axis tick labels},%
      axis background@ style={/pgfplots/on layer=axis background},%
      3d box foreground style={/pgfplots/on layer=axis foreground},%
    },
}

\newcommand{\etape}[2]{node (p#1) [etape] {#2}}

\newcommand{\transreceptor}[3]{%
  \path [linepart] (#1.east) -- node [above] {\scriptsize #2} (#3);}
\tikzstyle{matheq} = [node distance=8.75cm, text width=21em, minimum width=1cm,
minimum height=2em, text centered,color=black!50]

\renewcommand{\baselinestretch}{1.2}

%
\newcommand{\gkc}{\|\gk\|^{1/2}}
\newcommand{\xk}{x_k}
\newcommand{\Hk}{H_k}
\newcommand{\gk}{g_k}
\newcommand{\dk}{d_k}
\newcommand{\gkn}{g_{k+1}}
\newcommand{\xkn}{x_{k+1}}
\newcommand{\dkn}{d_{k+1}}
\newcommand{\lk}{\lambda_{k+1}}
\newcommand{\cmark}{\ding{51}}%
\newcommand{\xmark}{\ding{55}}%
\newcommand{\arca}{\texttt{CubicReg-A}}

\newtheorem{theorem}{Theorem}[section]
\newtheorem{lemma}[theorem]{Lemma}

\theoremstyle{definition}
\newtheorem{definition}[theorem]{Definition}
\newtheorem{example}[theorem]{Example}
\newtheorem{xca}[theorem]{Exercise}

\theoremstyle{remark}
\newtheorem{remark}[theorem]{Remark}
\newtheorem{corollary}{Corollary}[section]
% \newtheorem{lemma}{Lemma}[section]
\newtheorem{proposition}{Proposition}[section]
\newtheorem{condition}{Condition}[section]
% \newtheorem{definition}{Definition}[section]
% \newtheorem{example}{Example}[section]
\newtheorem{question}{Question}[section]
% \newtheorem{remark}{Remark}[section]
\newtheorem{assumption}{Assumption}[section]
\newtheorem{strategy}{Strategy}[section]


\def\algorithmautorefname{Algorithm}
\def\subsectionautorefname{Section}
\def\definitionautorefname{Definition}
\def\corollaryautorefname{Corollary}
\def\lemmaautorefname{Lemma}
\def\assumptionautorefname{Assumption}
\def\propositionautorefname{Proposition}
\def\tableautorefname{Table}
\def\strategyautorefname{Strategy}
\def\remarkautorefname{Remark}
\def\conditionautorefname{Condition}


\newcommand{\hsodm}{\texttt{HSODM}}
\newcommand{\hsodmarc}{\texttt{HSODMArC}}
\newcommand{\hsodmhvp}{\texttt{HSODM-HVP}}
\newcommand{\drsom}{\texttt{DRSOM}}
\newcommand{\drsomh}{\texttt{DRSOM-H}}
\newcommand{\lbfgs}{\texttt{LBFGS}}
\newcommand{\newtontr}{\texttt{Newton-TR}}
\newcommand{\cg}{\texttt{CG}}
\newcommand{\arc}{\texttt{ARC}}
\newcommand{\atrms}{\texttt{NO-ATR}}
\newcommand{\itrace}{\texttt{TRACE}}
\newcommand{\gd}{\texttt{GD}}
\newcommand{\utr}{\texttt{UTR}}
\newcommand{\utrhvp}{\texttt{UTR-HVP}}
\newcommand{\newtontrst}{\texttt{Newton-TR-STCG}}
\newcommand{\iutr}{\texttt{iUTR(Hessian)}}
\newcommand{\iutrhvp}{\texttt{iUTR}}

\newcommand{\atr}{\texttt{GL-ATR}}

\makeatletter
\renewcommand\paragraph{\@startsection{paragraph}{4}%
  \z@ \z@ {-.5em}%
  {\normalfont\normalsize\bfseries}}%
\makeatother


\usepackage{float}
\usepackage{etoolbox}
\usepackage{cleveref}

% Define subroutine float
\newfloat{subroutine}{htbp}{los}
\floatname{subroutine}{Subroutine}

% Define cref name for subroutine
\crefname{subroutine}{subroutine}{subroutines}
\Crefname{subroutine}{Subroutine}{Subroutines}

\allowdisplaybreaks